\notechapter{N-20220710}{Vectors, Ordered Fields, and the First Axioms of Real Numbers}

\section{Plane vectors}

The first page contains exercises on plane vectors.  A vector in the plane may be written
\[
  \vec a=(x,y)=x\mathbf i+y\mathbf j.
\]
Addition and scalar multiplication are coordinatewise:
\[
  (x,y)+(z,w)=(x+z,y+w),\qquad
  \lambda(x,y)=(\lambda x,\lambda y).
\]
The length and dot product are
\[
  |(x,y)|=\sqrt{x^2+y^2},\qquad
  (x,y)\cdot(z,w)=xz+yw.
\]
Thus the angle between two nonzero vectors satisfies
\[
  \cos\theta=\frac{\vec a\cdot\vec b}{|\vec a|\,|\vec b|}
  =\frac{xz+yw}{\sqrt{x^2+y^2}\sqrt{z^2+w^2}}.
\]
If
\[
  \vec x=(\cos\alpha,\sin\alpha),\qquad
  \vec y=(\cos\beta,\sin\beta),
\]
then
\[
  \cos\theta=\cos\alpha\cos\beta+\sin\alpha\sin\beta
  =\cos(\alpha-\beta).
\]
This gives a vector proof of the cosine subtraction formula.

\section{Field axioms}

The second page moves to the algebraic structure of the real numbers.  The field axioms are:
\begin{enumerate}[(F1)]
  \item \(x+(y+z)=(x+y)+z\);
  \item \(x+y=y+x\);
  \item there exists \(0\) such that \(0+x=x\);
  \item for each \(x\), there is \(-x\) such that \(x+(-x)=0\);
  \item \(x(yz)=(xy)z\);
  \item \(xy=yx\);
  \item there exists \(1\neq0\) such that \(1x=x\);
  \item for \(x\neq0\), there is \(x^{-1}\) such that \(xx^{-1}=1\);
  \item \(x(y+z)=xy+xz\).
\end{enumerate}

The note emphasizes that a ``space'' in mathematics is usually a set equipped with extra structure.  A field is not merely a set; it is a set equipped with addition, multiplication, identities, inverses, and compatibility laws.

\section{Ordered fields}

The order axioms include:
\begin{enumerate}[(O1)]
  \item Transitivity: if \(x\le y\) and \(y\le z\), then \(x\le z\).
  \item Antisymmetry: if \(x\le y\) and \(y\le x\), then \(x=y\).
  \item Totality: for all \(x,y\), either \(x\le y\) or \(y\le x\).
  \item Compatibility with addition: \(x\le y\Rightarrow x+z\le y+z\).
  \item Compatibility with multiplication: \(0\le x\) and \(0\le y\Rightarrow0\le xy\).
\end{enumerate}

The Archimedean axiom says that for every \(x>0\) and every \(y\), there exists a positive integer \(n\) such that
\[
  nx\ge y.
\]
A typical consequence is that every interval \((0,a)\) with \(a>0\) is nonempty: choose \(n\) so large that \(n>a^{-1}\), then \(1/n<a\).

\section{Elementary consequences of the order axioms}

The page records several exercises:
\begin{enumerate}[(i)]
  \item \(x\ge0\) iff \(-x\le0\), and \(y>1\) implies \(0<1/y<1\).
  \item \(1>0\) and \(-1\neq1\).
  \item If \(x\le y\) and \(a\le0\), then \(ax\ge ay\).
  \item If \(a\le b\) and \(x\le y\), then \(a+x\le b+y\), with equality iff \(a=b\) and \(x=y\).
  \item If \(a<x\) for every \(a<x\), then \(x\le y\); this is a typical contradiction argument using an interval between two distinct real numbers.
  \item \(x^2\ge0\) for every real \(x\).
  \item If \(a^2<a\), then \(0<a<1\).
  \item If nonzero \(x,y\) have the same sign, then
  \[
    (x+y)^2>(x-y)^2.
  \]
\end{enumerate}

\section[Embedding Q into R]{Embedding \(\mathbb Q\) into \(\mathbb R\)}

The note then proves that \(\mathbb R\) contains a copy of \(\mathbb Q\).  Define for integers \(n\)
\[
  \iota(n)=\underbrace{1+\cdots+1}_{n\text{ times}}
\]
for \(n>0\), \(\iota(0)=0\), and \(\iota(n)=-\iota(-n)\) for \(n<0\).  Then define
\[
  \iota\left(\frac pq\right)=\frac{\iota(p)}{\iota(q)}.
\]
One must check that this is well-defined: if \(p/q=s/t\), then \(pt=qs\), so
\[
  \iota(p)\iota(t)=\iota(q)\iota(s),
\]
and therefore
\[
  \frac{\iota(p)}{\iota(q)}=\frac{\iota(s)}{\iota(t)}.
\]
This embedding preserves addition, multiplication, and order.

\section{Nested intervals}

The final part introduces the nested interval axiom.  If
\[
  I_n=[a_n,b_n]
\]
is a descending sequence of closed intervals with real endpoints, then
\[
  I_1\supseteq I_2\supseteq I_3\supseteq\cdots
\]
implies
\[
  \bigcap_{n\ge1}I_n\neq\varnothing.
\]
Together with the field, order, and Archimedean axioms, this is one way to characterize the real numbers.


\section{Ordered fields and the real line}

An ordered field is a field equipped with an order compatible with addition and multiplication.  Compatibility means
\[
  a<b\Rightarrow a+c<b+c,\qquad
  0<a,\ 0<b\Rightarrow0<ab.
\]
The rational numbers form an ordered field, but they are not complete.  The real numbers are obtained by adding the missing limits, for example through Dedekind cuts or Cauchy sequences.

This is why early vector and coordinate exercises naturally lead to axioms of the real numbers.  Geometry needs coordinates, coordinates need arithmetic, and analysis needs completeness.

\section{Vectors as algebraicized geometry}

A plane vector \((x,y)\) packages a displacement into two real numbers.  Vector addition corresponds to composing displacements; scalar multiplication corresponds to stretching.  The algebraic laws of vectors are therefore not arbitrary axioms.  They express basic geometric operations in coordinate form.

\section{The conceptual payoff}

The note moves from computational vector exercises to the axiomatic construction of the real line.  The important conceptual shift is that familiar formulas are only the visible part of an axiomatic structure.  A rigorous real-number theory must explain why rational numbers embed, why intervals contain points, and why irrational numbers such as \(\sqrt2\) can exist.


\paragraph{Expanded synthesis.}
For this chapter, the elementary material should be read as a study of what remains invariant when coordinates change. The earlier sections (`Plane vectors`, `Field axioms`, `Ordered fields`, `Elementary consequences of the order axioms`, `Embedding \(\mathbb Q\) into \(\mathbb R\)`) compute with arrays, bases, pivots, determinants, or eigenvectors; the deeper object is a linear map together with the spaces it acts on. A matrix is a representation of that map after bases have been chosen. This is why formulas such as
\[
  A' = P^{-1}AP,\qquad \widetilde A = T^{-1}AS
\]
are not cosmetic: they distinguish an intrinsic morphism from a coordinate description.

The structural hierarchy is as follows. Kernels record what a map forgets, images record what it reaches, cokernels record obstructions to solvability, and quotients record information after an equivalence has been imposed. Determinants act on the top exterior power and therefore measure oriented volume; traces are invariant under cyclic rearrangement and become characters in representation theory; adjoints depend on an inner product and lead to self-adjoint spectral theory.

A useful way to return to the computations is to ask, for every formula, which invariant it is measuring. Row reduction measures rank and kernel; diagonalization measures decomposition into invariant subspaces; Gram--Schmidt measures orthogonal projection; positive definiteness measures ellipsoid geometry; tensor notation measures how components transform. The elementary methods are therefore the finite-dimensional laboratory in which duality, quotienting, functoriality, and spectral decomposition can be seen clearly.


\section{From coordinates to axioms}

This note jumps from vectors to field/order axioms.  The bridge is that coordinates only work because the scalars behave well.  When we write
\[
  (x,y)+(z,w)=(x+z,y+w),
\]
we are silently using the field structure of \(\mathbb R\).  When we compare lengths or prove inequalities, we are using the ordered-field structure.  When we take limits or talk about intervals containing points, we need completeness.

So the lesson is not merely that \(\mathbb R\) has many axioms.  It is that every familiar calculation depends on a layer of structure: algebra, order, and completeness.
